% ============================================
% SECTION 3: PROBLEM STATEMENT
% ============================================
\subsection{Problem Statement}

\begin{frame}{Problem Statement}
    \begin{alertblock}{Core Challenge}
        Accurately estimating a camera's geographic coordinates and viewing direction from a single mountain photograph when GNSS/GPS signals are unavailable.
    \end{alertblock}

    \vspace{1em}

    \textbf{Current Technical Barriers:}
    \begin{enumerate}
        \item \textbf{Visual Obstruction:} Haze, clouds, and low contrast frequently obscure the skyline boundary.
        \item \textbf{Search Scale:} Matching a partial skyline against a massive database without knowing the camera's initial heading is computationally expensive.
        \item \textbf{Data Scarcity:} A critical lack of annotated mountain datasets for the Himalayan region compared to European or Chinese terrain.
    \end{enumerate}
\end{frame}

\begin{frame}{Problem Analysis}
    \textbf{Impact of GPS Failure in Mountains:}
    \begin{itemize}
        \item \textbf{Topographic Shielding:} Deep Himalayan valleys block line-of-sight to satellites.
        \item \textbf{Safety Risk:} Navigation failure in the Khumbu region can be fatal for trekkers and rescue teams.
        \item \textbf{Infrastructure Limits:} Remote areas lack cell towers or internet connectivity for cloud-based maps.
    \end{itemize}

    \vspace{1em}

    \begin{exampleblock}{Significance of the Solution}
        Mountain skylines are unique, permanent, and visible even when satellites are not. This project provides a robust, offline navigation alternative using the observer's own camera as a sensor.
    \end{exampleblock}
\end{frame}
